<<<<<<< HEAD
\chapter{Results}{\label{ch:results}}

=======
\chapter{Results}
\label{chap:results}

This chapter reports the evaluation of each component of HepatoTwin, in the
order set out in Section~\ref{sec:m_eval}. Section~\ref{sec:r_layer2} covers
the liver-state model, its uncertainty and the fibrosis triage;
Section~\ref{sec:r_fairness} the fairness audit; Section~\ref{sec:r_nhanes3}
the external and long-term validation on NHANES~III;
Sections~\ref{sec:r_priorities}--\ref{sec:r_summary} the evidence engine, the
recipe ranking engine and the plain-language summaries; and
Sections~\ref{sec:r_calibration}--\ref{sec:r_insilico} the calibration of the
simulator and the in-silico evaluation of the updating twin.
Section~\ref{sec:r_overview} brings the results together. Unless stated,
results are on held-out participants inside the model's scope (ages 12--59,
reliable examination), intervals are 95\% bootstrap intervals, and figures are
unweighted.

% ======================================================================
\section{Liver-State Model, Uncertainty and Fibrosis Triage}
\label{sec:r_layer2}

\subsection{Discrimination and calibration}

The development of the model is reported in Chapter~\ref{chap:experiments};
Table~\ref{tab:r_layer2} summarises the selected model on the 1{,}236 in-scope
test participants. Both classifiers rank participants well (AUC 0.881 for CAP
$\geq$ 248 and 0.879 for CAP $\geq$ 280) and are well calibrated: the expected
calibration error is 0.018 and 0.022, so a stated probability of 30\% means
about 30\% of such patients have that much liver fat. On the 1{,}100
participants with every input to both scores, the CAP $\geq$ 248 classifier
outperformed the Fatty Liver Index (0.861) and the Hepatic Steatosis Index
(0.836), with paired differences that exclude zero
(Chapter~\ref{chap:experiments}).

\begin{table}[htbp]
\centering
\caption{Selected liver-state model on the in-scope test participants ($n = 1{,}236$).}
\label{tab:r_layer2}
\begin{tabular}{p{6.2cm}p{3.6cm}p{3.6cm}}
\toprule
Measure & CAP $\geq$ 248~dB/m & CAP $\geq$ 280~dB/m \\
\midrule
Prevalence & 0.506 & 0.340 \\
ROC-AUC (95\% CI) & 0.881 (0.863--0.899) & 0.879 (0.860--0.898) \\
Brier score & 0.140 & 0.132 \\
Expected calibration error & 0.018 & 0.022 \\
F1 at the decision threshold & 0.806 (threshold 0.432) & 0.713 (threshold 0.365) \\
\midrule
CAP $R^2$ (95\% CI) & \multicolumn{2}{l}{0.562 (0.517--0.601)} \\
CAP RMSE (dB/m) & \multicolumn{2}{l}{41.4 (39.3--43.6)} \\
\bottomrule
\end{tabular}
\end{table}

\subsection{Uncertainty}

The 90\% conformal interval for CAP is the estimate $\pm$66.6~dB/m. It
contained the measured CAP for 90.9\% of in-scope test participants (target
90\%): 93.8\% at ages 12--17, 90.8\% at 18--39 and 89.7\% at 40--59. The 80\%
and 95\% intervals covered 80.1\% and 95.6\%. The interval is wide---it spans
more than the distance between the S1 and S3 cut-offs---which is why patients
are shown the calibrated probabilities instead of a CAP estimate. The
bootstrap bands around the calibrated probabilities had a mean width of 0.054
(CAP $\geq$ 248) and 0.051 (CAP $\geq$ 280), so the calibration step adds
little uncertainty.

\subsection{Fibrosis triage}

Table~\ref{tab:r_fib4} compares FIB-4 with liver stiffness measured in the
same FibroScan examination, for the 7{,}586 in-scope adults with a reliable
reading. FIB-4 discriminates poorly in this young, community population (AUC
0.59 for stiffness $\geq$ 8~kPa and 0.67 for $\geq$ 12~kPa), and its low
cut-off misses most raised stiffness (sensitivity 20\% at 8~kPa). Its value is
in ruling out: 92\% of adults with FIB-4 below 1.3 had stiffness below
8~kPa, and 98\% below 12~kPa. The high cut-off is specific (99.6\%) but rare:
0.9\% of adults were classed as high risk and 8.3\% as indeterminate. This is
why the twin uses a reliable measured stiffness in place of FIB-4 whenever one
is available, and why the summary tells young adults with probable liver fat
that a low FIB-4 does not fully rule out scarring.

\begin{table}[htbp]
\centering
\caption{FIB-4 against measured liver stiffness, in-scope adults with a reliable reading ($n = 7{,}586$; test partition alone, $n = 1{,}513$, in brackets).}
\label{tab:r_fib4}
\begin{tabular}{p{3.6cm}p{1.9cm}p{1.9cm}p{1.9cm}p{1.9cm}p{1.9cm}}
\toprule
Stiffness, rule & AUC & Sensitivity & Specificity & PPV & NPV \\
\midrule
$\geq$ 8~kPa, FIB-4 $\geq$ 1.3 & 0.59 (0.62) & 0.20 & 0.92 & 0.19 & 0.92 \\
$\geq$ 8~kPa, FIB-4 $>$ 2.67 & & 0.05 & 0.996 & 0.54 & 0.92 \\
$\geq$ 12~kPa, FIB-4 $\geq$ 1.3 & 0.67 (0.72) & 0.33 & 0.92 & 0.11 & 0.98 \\
$\geq$ 12~kPa, FIB-4 $>$ 2.67 & & 0.12 & 0.994 & 0.39 & 0.97 \\
\bottomrule
\end{tabular}
\end{table}

% ======================================================================
\section{Fairness Audit}
\label{sec:r_fairness}

Table~\ref{tab:r_fairness} reports the CAP $\geq$ 280 classifier by sex, race
and ethnicity, and age. Discrimination is similar in every group: AUC ranges
from 0.840 to 0.894 across sexes and ethnic groups, and every interval
overlaps every other. It is highest at ages 12--17 (0.931) and lowest at
40--59 (0.841), where liver fat is most common and harder to separate.

Calibration differs by ethnicity. The mean probability was too high for
non-Hispanic Black participants (0.337 against an observed share of 0.290; the
interval of the difference, 0.002 to 0.093, excludes zero) and too low for
non-Hispanic Asian participants (0.251 against 0.313; interval $-0.129$ to
0.002, borderline). The pattern for Asian participants agrees with the known
tendency to more liver fat at a lower BMI. The CAP $\geq$ 248 classifier showed
the same direction for both groups, with intervals that included zero.

\begin{table}[htbp]
\centering
\caption{CAP $\geq$ 280 classifier by group, in-scope test participants.}
\label{tab:r_fairness}
\begin{tabular}{p{4.2cm}rp{3cm}rr}
\toprule
Group & $n$ & AUC (95\% CI) & Mean probability & Observed share \\
\midrule
All in scope & 1{,}236 & 0.879 (0.859--0.896) & 0.336 & 0.340 \\
Men & 583 & 0.889 (0.861--0.915) & 0.379 & 0.376 \\
Women & 653 & 0.869 (0.838--0.894) & 0.299 & 0.308 \\
Mexican American & 174 & 0.884 (0.831--0.930) & 0.435 & 0.477 \\
Other Hispanic & 141 & 0.869 (0.809--0.921) & 0.291 & 0.291 \\
Non-Hispanic White & 461 & 0.894 (0.866--0.920) & 0.337 & 0.343 \\
Non-Hispanic Black & 241 & 0.876 (0.828--0.920) & 0.337 & 0.290 \\
Non-Hispanic Asian & 134 & 0.840 (0.768--0.900) & 0.251 & 0.313 \\
Other or multiracial & 85 & 0.859 (0.759--0.936) & 0.337 & 0.306 \\
Ages 12--17 & 209 & 0.931 (0.887--0.969) & 0.142 & 0.144 \\
Ages 18--39 & 512 & 0.872 (0.841--0.902) & 0.290 & 0.285 \\
Ages 40--59 & 515 & 0.841 (0.804--0.875) & 0.461 & 0.474 \\
\bottomrule
\end{tabular}
\end{table}

The recipe ranking treated groups alike. The share of top-ten recipes meeting
all three criteria was 0.810 for men and 0.839 for women, between 0.800 and
0.846 across the three ethnic groups with at least 30 patients, and between
0.793 and 0.827 across age groups, with overlapping intervals; every top-ten
recipe met every health reference in every group.

% ======================================================================
\section{External and Long-Term Validation on NHANES III}
\label{sec:r_nhanes3}

\subsection{Discrimination against ultrasound}

On 10{,}682 adults examined 30 years before the training data, with a
different reference standard and 30 of the 81 inputs missing, the model ranked
moderate-to-severe steatosis on ultrasound better than the established scores
computed from the same records (Table~\ref{tab:r_nhanes3}). Its AUC of 0.771 is
lower than the 0.879 achieved against CAP in the training survey, as expected
when the reference standard, the era and the available inputs all change, but
it is above the Hepatic Steatosis Index (0.725) and, on the participants with
GGT, the Fatty Liver Index (0.744 against 0.766 for the model). The AUC was
0.783 in men and 0.755 in women, 0.787 in Mexican American and 0.790 in
non-Hispanic White participants, and 0.690 in non-Hispanic Black participants,
in whom moderate-to-severe steatosis was also least common (16\%). The mean
probability (0.255) was close to the ultrasound share (0.207), and the share
rose in every tenth of the probability from the fourth upwards.

\begin{table}[htbp]
\centering
\caption{Moderate-to-severe steatosis on ultrasound, NHANES~III adults aged 20--59.}
\label{tab:r_nhanes3}
\begin{tabular}{p{7cm}rp{3.5cm}}
\toprule
Score & $n$ & AUC (95\% CI) \\
\midrule
HepatoTwin: probability of CAP $\geq$ 280 & 10{,}682 & 0.771 (0.759--0.782) \\
Hepatic Steatosis Index & 10{,}055 & 0.725 (0.712--0.738) \\
Waist circumference & 10{,}374 & 0.716 (0.703--0.728) \\
BMI & 10{,}668 & 0.698 (0.686--0.712) \\
\midrule
\emph{Participants with GGT:} Fatty Liver Index & 7{,}649 & 0.744 (0.730--0.758) \\
\emph{Participants with GGT:} HepatoTwin & 7{,}649 & 0.766 (0.752--0.780) \\
\bottomrule
\end{tabular}
\end{table}

\subsection{Long-term mortality}

Over a median of 27.3 years, 2{,}173 of 10{,}671 linked participants died. The
death rate rose with the model's probability: 4.0, 6.8, 7.9 and 13.8 deaths
per 1{,}000 person-years from the lowest to the highest quarter.
Table~\ref{tab:r_mortality} gives rate ratios adjusted for age and sex.
Participants in the highest quarter had a 58\% higher all-cause death rate and
more than twice the cardiovascular death rate of those in the lowest; cancer
deaths did not differ. The model's probability separated cardiovascular
mortality more strongly than the ultrasound grade itself, plausibly because it
also carries metabolic risk from waist, glucose and triglycerides. A high FIB-4
was associated with three times the all-cause death rate. These ratios are
adjusted for age and sex only: smoking, alcohol and income were not mapped,
and the public mortality file has no liver-related cause of death.

\begin{table}[htbp]
\centering
\caption{Age- and sex-adjusted death rate ratios, NHANES~III (median follow-up 27.3 years).}
\label{tab:r_mortality}
\begin{tabular}{p{2.6cm}p{3.6cm}p{3.6cm}p{3.6cm}}
\toprule
Deaths & Highest vs lowest quarter of the probability & Ultrasound moderate--severe vs not & FIB-4 high vs low \\
\midrule
All-cause & 1.58 (1.38--1.81) & 1.29 (1.17--1.41) & 3.23 (2.47--4.23) \\
Cardiovascular & 2.15 (1.64--2.81) & 1.35 (1.13--1.61) & 1.89 (1.01--3.54) \\
Cancer & 1.00 (0.79--1.28) & 1.05 (0.88--1.26) & 2.63 (1.57--4.41) \\
\bottomrule
\end{tabular}
\end{table}

% ======================================================================
\section{Evidence Engine}
\label{sec:r_priorities}

Among the 300 test participants of the ranking evaluation, the evidence engine
produced 44 distinct priority profiles when the 24-hour recall was used, and
the most common profile covered 8\% of patients. Without the recall, the same
patients fell into 9 profiles and the most common covered 43\%: the patient's
own intake is what makes the priorities personal.

The priority levels change the recipes in the intended direction
(Table~\ref{tab:r_response}). With a high saturated-fat priority the top ten
averaged 3.87\% of energy from saturated fat, against 4.86\% with a low one;
with a high sugar priority, 4.39\% from sugars against 4.62\%. Fibre did not
follow the level, because almost every recommended recipe already exceeds the
fibre reference by more than the margin counts.

\begin{table}[htbp]
\centering
\caption{Mean composition of the top-ten recipes by priority level (300 patients).}
\label{tab:r_response}
\begin{tabular}{p{5cm}rrr}
\toprule
Measure (priority) & Low & Medium & High \\
\midrule
Saturated fat, \% of energy (limit saturated fat) & 4.86 & 3.78 & 3.87 \\
Sugars, \% of energy (limit excess sugar) & 4.62 & 4.37 & 4.39 \\
Fibre, g per 1{,}000~kcal (diet quality) & 24.9 & 35.3 & 32.3 \\
\bottomrule
\end{tabular}
\end{table}

% ======================================================================
\section{Recipe Evaluation and Ranking Engine}
\label{sec:r_ranking}

\subsection{Comparison with alternative rankings}

Table~\ref{tab:r_methods} compares the proposed ranking with the alternatives
on 300 held-out patients. The health-gated ranking with personal fit met all
three criteria at once for 82.3\% of top-ten recipes, and every recommended
recipe met every one of the patient's health references. Weighted scoring
without the gate reached a higher share for cuisine and feasibility but met
the health references for only 27.8\% of recipes, so only 27.2\% met all
three: without the gate, convenient and familiar recipes displace healthy
ones. Ranking on health alone met the references but rarely the patient's
cuisine or the practical limits (2.2\%), and random selection almost never met
all three (0.3\%). Diversity re-ranking raised the diversity of the list only
slightly (0.963 against 0.949) at a small cost in practicality, so it is not
the default.

\begin{table}[htbp]
\centering
\caption{Share of top-ten recipes meeting each criterion, 300 held-out patients.}
\label{tab:r_methods}
\begin{tabular}{p{6.4cm}rrrrr}
\toprule
Method & Health & Cuisine & Feasibility & All three & Diversity \\
\midrule
Health-gated with personal fit (proposed) & 1.000 & 0.976 & 0.832 & \textbf{0.823} & 0.949 \\
Health-gated without personal fit & 1.000 & 0.975 & 0.831 & 0.826 & 0.952 \\
Health-gated with diversity re-ranking & 1.000 & 0.972 & 0.820 & 0.804 & 0.963 \\
Health-gated without acceptability & 1.000 & 0.980 & 0.807 & 0.803 & 0.945 \\
Weighted scoring, no gate & 0.278 & 0.998 & 0.957 & 0.272 & 0.948 \\
Health score only & 1.000 & 0.068 & 0.300 & 0.022 & 0.963 \\
Nutrition only (not personalised) & 1.000 & 0.074 & 0.300 & 0.022 & 0.968 \\
Random & 0.057 & 0.053 & 0.596 & 0.003 & 0.962 \\
\bottomrule
\end{tabular}
\end{table}

Across 36 combinations of the health, cuisine and feasibility weights, the
median overlap of the top ten with the default was 0.78; it fell most when the
cuisine weight was reduced to 0.1. The ranking is therefore not fragile to the
exact weights, except that cuisine needs a meaningful share.

\subsection{Personalisation}

A review of sample summaries by the project team found that patients without
food answers were offered almost the same dishes. Before personal fit, two
patients of the same cuisine shared 93\% of their top-ten recipes, and 300
patients received only 28 distinct lists (Table~\ref{tab:r_personal}). Adding
personal fit, applied only to recipes that pass the health gate, made the
lists clearly personal---124 distinct lists and 71\% overlap at weight
0.4---while every recipe still met every health reference and the share
meeting all three criteria barely changed (82.3\% against 82.6\%). The 120
patients without a reliable recall keep the gated order, as nothing is known
about their food; in use, every patient also gives favourite dishes and liked
ingredients at sign-up.

\begin{table}[htbp]
\centering
\caption{Personalisation of the top ten, 300 held-out patients.}
\label{tab:r_personal}
\begin{tabular}{p{5.2cm}rrrr}
\toprule
 & Without recall & With recall & Personal fit 0.2 & Personal fit 0.4 \\
\midrule
Distinct top-ten lists & 15 & 28 & 81 & \textbf{124} \\
Overlap between patients of the same cuisine & 0.95 & 0.93 & 0.85 & \textbf{0.71} \\
Top ten meeting all three criteria & -- & 0.826 & 0.834 & 0.823 \\
\bottomrule
\end{tabular}
\end{table}

\subsection{Weekly meal plans}

Every one of the 300 seven-day plans was complete and met every weekly
constraint: every meal and every day met the patient's references, no recipe
was repeated beyond the limits, the two meals of each day had different
protein families (six families a week on average), and every meal was in the
patient's cuisine or a declared similar one. The requested cuisine alone
filled 58\% of weeks; allowing a favourite to appear twice raised this to 80\%,
and borrowing from similar cuisines completed the rest, with 0.2 borrowed
meals per plan on average. Borrowing was needed almost only for Canadian
recipes (98\% of those weeks), the smallest cuisine with too few qualifying
meals. On average 81\% of planned meals met the practical limits.

% ======================================================================
\section{Plain-Language Summaries}
\label{sec:r_summary}

Every one of 300 summaries was faithful: each number and evidence citation
appeared in the twin's facts (Table~\ref{tab:r_summaries}). The check caught
every planted error that changed or added a number or invented a citation, and
never flagged a harmless rewording. It cannot catch a wrong word that carries
no number (``low'' written for ``high''), which it never caught; such errors
are prevented by writing the summary from a template, and would need a
separate check if a language model wrote the text.

\begin{table}[htbp]
\centering
\caption{Plain-language summaries, 300 in-scope test patients, before and after the team's review.}
\label{tab:r_summaries}
\begin{tabular}{p{7.4cm}p{3cm}p{3cm}}
\toprule
Check & Before & After \\
\midrule
Faithful summaries & 100\% & 100\% \\
Planted number changed, number added, citation invented: caught & 100\% & 100\% \\
Planted harmless rewording: flagged & 0\% & 0\% \\
Planted wrong word with no number: caught & 0\% & 0\% \\
Flesch Reading Ease, median & 43 & 56.5 \\
Flesch--Kincaid grade, median & 12.5 & 9.8 \\
Evidence codes shown to the patient & Yes & No \\
\bottomrule
\end{tabular}
\end{table}

The team's review of ten summaries found that every statement matched the
twin's output, but that the text was hard for a patient to use: evidence codes
meant nothing to a lay reader, liver scarring and glycaemic load were not
explained, scores were given without the published ranges, and two notices
were shown where they did not apply (a teenager without liver fat told to see
a doctor before changing diet, and a young adult without liver fat advised to
ask about elastography). After the changes described in
Section~\ref{sec:m_summary} the median reading grade fell from 12.5 to 9.8,
while faithfulness was unchanged.

% ======================================================================
\section{Simulator Calibration}
\label{sec:r_calibration}

Table~\ref{tab:r_calibration} compares the tuned simulator with the published
outcomes. The calibration targets were reproduced closely: Look AHEAD's
one-year mean loss ($-8.4\%$ against $-8.6\%$), its spread and the shares
losing 5\% and 10\%; the DPP's share reaching the 7\% goal at 24 weeks (49.9\%
against 50\%); and the shape of the weight course and the weight-loss groups
of Vilar-Gomez et al. The held-back targets show where the simulator is
weaker: virtual Look AHEAD participants had kept off less weight at eight
years than the real ones ($-1.9\%$ against $-4.7\%$), and DPP participants less
at about three years (3.3~kg against 5.6~kg). The simulator therefore
represents long-term maintenance pessimistically. For the twelve months
evaluated here the first-year targets matter most, and these were met.

\begin{table}[htbp]
\centering
\caption{Simulator outcomes against published trial results.}
\label{tab:r_calibration}
\begin{tabular}{p{7.6cm}p{1.6cm}rr}
\toprule
Outcome & Use & Published & Simulated \\
\midrule
Look AHEAD year 1, mean change, intensive (\% of weight) & Tuned & $-8.6$ & $-8.4$ \\
Look AHEAD year 1, mean change, usual care (\%) & Tuned & $-0.6$ & $+0.1$ \\
Look AHEAD year 1, lost $\geq$ 5\% / $\geq$ 10\%, intensive (\%) & Tuned & 68.0 / 37.7 & 66.9 / 39.7 \\
DPP, reached the 7\% goal at 24 weeks (\%) & Tuned & 50.0 & 49.9 \\
Vilar-Gomez, mean change at 52 weeks (\%) & Tuned & $-3.8$ & $-3.0$ \\
Vilar-Gomez, lost $<$5 / 5--7 / 7--10 / $\geq$10\% (\%) & Tuned & 70 / 12 / 9 / 10 & 66 / 11 / 11 / 12 \\
Vilar-Gomez, lost $\geq$ 5\% (\%) & Check & 30.0 & 33.8 \\
Look AHEAD year-1 regain, third regaining most (points) & Selection & 1.5 & 2.0 \\
\midrule
Look AHEAD year 8, mean change, intensive (\%) & Held back & $-4.7$ & $-1.9$ \\
Look AHEAD year 8, lost $\geq$ 5\% / $\geq$ 10\% (\%) & Held back & 50.3 / 26.9 & 40.9 / 29.4 \\
DPP, met the goal at the last visit (\%) & Held back & 38.0 & 33.7 \\
DPP, mean loss at about 3 years (kg) & Held back & 5.6 & 3.3 \\
\bottomrule
\end{tabular}
\end{table}

With the adherence of routine advice to patients with fatty liver, the virtual
cohort lost 1.7\% of its weight by month 3 on Scenario~A ($-500$~kcal a day)
and 0.8\% by month 12, as most patients eased off; at month 12, 25.7\% were at
least 5\% below their starting weight, against 15.6\% on usual diet, whose mean
weight drifted up by 1.3\%. Mean CAP fell by 3.7~dB/m on Scenario~A and 2.2~dB/m on Scenario~B.

% ======================================================================
\section{In-Silico Evaluation of the Updating Twin}
\label{sec:r_insilico}

\subsection{Weight forecasts}

Table~\ref{tab:r_forecasts} reports the twin's weight forecasts for adults on
the two diet scenarios. The twin's error was 1.9~kg three months ahead and
2.8~kg six months ahead, against 2.9 and 4.7~kg for the textbook curve that
does not learn, and its 90\% ranges contained the true weight 91\% of the
time, as a 90\% range should. Carrying the last weight forward was as accurate
as the twin (1.9 and 2.7~kg), because in the calibrated simulator most
patients ease off within months and their weight flattens; a flat forecast is
hard to beat when weight truly is flat. Unlike the last weight, the twin also
gives a calibrated range and a projection under each scenario. Without the
patient's answers about what they followed, the six-month error rose from 2.77
to 3.03~kg, so the check-in answers add information.

\begin{table}[htbp]
\centering
\caption{Weight forecast error for adults (mean absolute error, kg) and coverage of the 90\% range.}
\label{tab:r_forecasts}
\begin{tabular}{p{3.4cm}rrrrrr}
\toprule
Forecast & $n$ & Twin & Textbook & Last weight & Twin coverage & Twin, no answers \\
\midrule
Month 3 $\to$ 6 & 1{,}996 & 1.92 & 2.89 & 1.88 & 0.916 & 2.08 \\
Month 6 $\to$ 9 & 1{,}916 & 1.97 & 3.91 & 1.97 & 0.907 & 2.09 \\
Month 6 $\to$ 12 & 1{,}916 & 2.77 & 4.70 & 2.71 & 0.907 & 3.03 \\
\bottomrule
\end{tabular}
\end{table}

\textbf{Learning.} The twin's three-month-ahead error stayed at about 1.9~kg
from month 1 (1.98~kg) to month 9 (1.88~kg), while the textbook curve's error
grew from 2.25 to 4.75~kg as patients drifted away from the textbook path; the
twin's 90\% coverage stayed between 91\% and 93\%. Learning each patient's
course therefore keeps the forecast accurate as the textbook assumption
becomes less true, rather than reducing an already small error.

\textbf{Robustness.} The six-month forecast error was 2.77~kg with honest
answers, 2.81~kg with 30\% more missed visits, 2.91~kg with a scale reading
within about 1~kg, 2.77~kg with answers inflated by each person's
under-reporting, and 3.03~kg with no answers; coverage stayed between 0.89 and
0.91. The twin degrades gracefully with noisier, biased or missing inputs.

\subsection{Causes of slow progress and clinician alerts}

When the twin judged a patient behind plan and named the plan as not followed,
the patient's true adherence was below 0.4 every time (12{,}811 check-ins).
When it named no progress despite effort, true adherence was 0.6 or more every
time with honest answers, and 83.5\% of the time with over-reported answers.

The twin could not, however, learn whether a patient's body responds slowly.
The learned body response correlated only 0.03 to 0.06 with the true one:
within twelve monthly weights, ordinary weight drift (about 1.3~kg per
square-root month) is larger than the true difference between bodies (about
20\%). An alert designed to flag slow bodies found them no better than chance
(precision 0.26 against a base rate of 0.20; recall 0.02). This result changed
the product: the alert now reports only what the twin observes---no progress
at two check-ins despite reported effort---and lists what the clinician should
check. Judged on that claim, 78\% of alerted patients were still behind plan at
month 12, against 63\% of all patients. Behind the alerts lay ordinary weight
drift (74\%) and a slow body (26\%), and, when answers were over-reported, also
intake higher than reported (8\%), which the suggested food diary would reveal.

\subsection{Detecting change}

A patient who was truly behind plan was flagged at the same check-in in 96.5\%
of cases and within a month in 99.4\%; 51 of 1{,}659 were never flagged, and
no patient was flagged without truly being behind. Because the observed weight
is the true weight in this simulation, detection here is limited only by the
check-in schedule and missed visits; real scales add noise. Liver fat falling below 280~dB/m was harder to detect. Of 277
patients on Scenario~A who started above it and were flagged at sign-up, 51
truly fell below it within the year; the twin's probability fell below 0.5 for
29 of them (median one month before the truth), but also for 46 patients whose
CAP stayed above 280. Close to the threshold, a calibrated probability cannot
place a person on one side.

\subsection{Tracking liver fat}

The probability of CAP $\geq$ 280, recomputed from the record the twin holds
at each time, stayed accurate and calibrated over the year
(Table~\ref{tab:r_liver}): its AUC was 0.88 at months 0, 6 and 12 on every
scenario, and its mean stayed close to the true share as liver fat fell (0.31
against 0.29 at month 12 on Scenario~A). At month 6, the twin's forecast of
month 12 was better than reusing the sign-up probability (AUC 0.875 against
0.868, Brier score 0.126 against 0.133) and its mean was right (0.294 against
a true share of 0.291, where the sign-up probability gave 0.337). Across
patients, the change in the probability from month 0 to month 12 followed the
true change in CAP (Spearman correlation 0.79). Without follow-up blood tests,
the month-12 AUC was unchanged (0.88), because in the simulator liver fat
changes through weight and waist.

\begin{table}[htbp]
\centering
\caption{Probability of CAP $\geq$ 280 against the simulated truth, Scenario~A ($n = 1{,}236$).}
\label{tab:r_liver}
\begin{tabular}{p{6.2cm}rrrr}
\toprule
Time & True share & Mean probability & AUC & Brier \\
\midrule
Month 0 & 0.338 & 0.336 & 0.879 & 0.132 \\
Month 6 & 0.294 & 0.302 & 0.891 & 0.119 \\
Month 12 & 0.290 & 0.310 & 0.880 & 0.124 \\
Month 12, forecast at month 6 ($n = 1{,}149$) & 0.291 & 0.294 & 0.875 & 0.126 \\
Month 12, sign-up probability reused ($n = 1{,}149$) & 0.291 & 0.337 & 0.868 & 0.133 \\
\bottomrule
\end{tabular}
\end{table}

\subsection{Ordering the scenarios}

At sign-up the twin projects Scenario~A below Scenario~B below usual diet for
every patient, and the simulated truth agreed for 95.6\% of adults. At month
6, its forecasts of month-12 weight ordered Scenario~A and Scenario~B as the
truth did for 95.6\% of patients. For liver fat the month-6 forecasts agreed
with the true order for only 54\% of patients: the true difference in CAP
between a 500 and a 250~kcal deficit after a year is small next to the
variation between people, so the twin can rank the scenarios on average but not
reliably for each person.

% ======================================================================
\section{Overview}
\label{sec:r_overview}

Table~\ref{tab:r_overview} brings the main results together. The liver-state
model is accurate and calibrated on held-out patients, carries to a cohort 30
years older with a different reference standard, and its probability
predicts long-term cardiovascular death. The ranking engine is the only method
whose recommendations always met the patient's health references, and personal
fit made its lists personal without giving up any health rule. In simulation,
the updating twin forecasts weight better than the textbook curve with
calibrated ranges, notices when a patient falls behind, tracks liver fat over a
year and degrades gracefully with poor data. The evaluation also exposed limits
that shaped the system: a body's response to diet cannot be identified within a
year, so the twin reports lack of progress rather than its cause; FIB-4 is weak
in young community patients, so measured stiffness is preferred; the model is
less well calibrated for some ethnic groups; and the first summaries were
faithful but hard to read.

\begin{table}[htbp]
\centering
\caption{Main results.}
\label{tab:r_overview}
\begin{tabular}{p{4.2cm}p{9.8cm}}
\toprule
Component & Result \\
\midrule
Liver-state model & AUC 0.881 / 0.879 (CAP $\geq$ 248 / 280), calibration error $\leq$ 0.022; above FLI and HSI \\
Uncertainty & 90\% conformal interval covered 90.9\% \\
Fairness & AUC 0.84--0.93 in every group; probability too high for non-Hispanic Black participants \\
NHANES III & AUC 0.771 against ultrasound (HSI 0.725, FLI 0.744); highest quarter: 1.58$\times$ all-cause and 2.15$\times$ cardiovascular death \\
Recipe ranking & 82.3\% of top-ten recipes meet all three criteria, 100\% meet every health reference; weighted scoring without the gate 27.2\% \\
Personalisation & 124 distinct top-ten lists (28 without personal fit); same-cuisine overlap 71\% (93\%) \\
Weekly plans & 100\% complete and meeting every constraint \\
Summaries & 100\% faithful; reading grade 12.5 $\to$ 9.8 after review \\
Twin forecasts (in silico) & 1.9 / 2.8~kg error three / six months ahead (textbook 2.9 / 4.7); 91\% coverage \\
Twin causes and alerts & ``Plan not followed'' always correct; body response not identifiable (correlation 0.03--0.06) \\
Liver tracking (in silico) & AUC 0.88 at months 0, 6 and 12; month-6 forecast beats the sign-up probability \\
\bottomrule
\end{tabular}
\end{table}
>>>>>>> c07f45ec22f6707d048f989b3f56b5180ae12beb
